\documentclass[11pt]{article}
\usepackage{xspace,mathrsfs}
\usepackage{graphicx,rotating}
\usepackage{url}
\usepackage{amscd,amsfonts,amsmath,amssymb}
\usepackage{color}
\usepackage{longtable}
\usepackage{multirow,setspace}
\renewcommand{\topfraction}{0.9}
\renewcommand{\textfraction}{0.1}
\usepackage{float}
\usepackage{booktabs}
\usepackage{caption}
\usepackage{fancyhdr}

\def\shownotes{0}   % set 1 for version with author notes
                    % set 0 for no notes

%%%%%%%%%%%%%%%%%%%%%%%%%%%%%%%%%%%%%%%%%%%%%%%%%%%%%%%%%%%%%%
% Some macros (you can ignore everything until "end of macros")
\topmargin 0pt \advance \topmargin by -\headheight \advance
\topmargin by -\headsep

\textheight 8.9in
\oddsidemargin 0pt \evensidemargin \oddsidemargin \marginparwidth
0.5in

\textwidth 6.5in

%%%%%%
\DeclareMathAlphabet{\mathsl}{OT1}{cmr}{m}{sl}
\DeclareMathAlphabet{\mathsc}{OT1}{cmr}{m}{sc}
\DeclareMathAlphabet{\mathslbf}{OT1}{cmr}{bx}{sl}

\newcommand{\procfont}[1]{\mathsc{#1}}
\newcommand{\algfont}[1]{{\mathsf{#1}}}
\newcommand{\schemefont}[1]{{\mathsf{#1}}}

\newcommand{\gamefont}[1]{\mathrm{#1}}
\newcommand{\getsr}{{{\,\leftarrow{\hspace*{-3pt}\raisebox{.75pt}{$\scriptscriptstyle\$$}}\,}}}
\newcommand{\bits}{\{0,1\}}
\newcommand{\xor}{\oplus}
\newcommand{\Adv}{\mathbf{Adv}}
\newcommand{\AdvGen}[3]{\mathbf{Adv}^{\mathrm{#1}}_{#2}(#3)}
\newcommand{\Prob}[1]{{\Pr\left[\,{#1}\,\right]}}
\newcommand{\main}{\procfont{Game}}
\newcommand{\gameName}[1]{\underline{$\main$ #1} \\[2pt]}
\newcommand{\return}{\textrm{Ret}~}

\newcommand{\gamesfontsize}{\small}
\renewcommand{\captionfont}{}
\newcommand{\indd}{\hspace*{12pt}}

\newcommand{\mpage}[2]{\begin{minipage}{#1\textwidth} #2 \end{minipage}}
\newcommand{\fpage}[2]{\framebox{\begin{minipage}[t]{#1\textwidth}\setstretch{1.2}\gamesfontsize  #2 \end{minipage}}}
\newcommand{\hfpages}[3]{\hfpagess{#1}{#1}{#2}{#3}}
\newcommand{\hfpagess}[4]{
        \begin{tabular}{c@{\hspace*{.5em}}c}
        \framebox{\begin{minipage}[t]{#1\textwidth}\setstretch{1.2}\gamesfontsize #3 \end{minipage}}
        &
        \framebox{\begin{minipage}[t]{#2\textwidth}\setstretch{1.2}\gamesfontsize #4 \end{minipage}}
        \end{tabular}
    }

\newcommand{\Initialize}{\procfont{Initialize}}
\newcommand{\Finalize}{\procfont{Finalize}}
\newcommand{\LR}{\procfont{LR}}
\newcommand{\Enc}{\procfont{Enc}}
\newcommand{\Vf}{\procfont{Vf}}
\newcommand{\InitializeCode}{\underline{\textrm{proc} $\Initialize$}\\[2pt]}
\newcommand{\FinalizeCode}[1]{\underline{\textrm{proc} $\Finalize(#1)$}\\[2pt]}
\newcommand{\FinalizeCodeNoArg}{\underline{\textrm{proc} $\Finalize$}\\[2pt]}
\newcommand{\LRCode}{\underline{\textrm{proc} $\LR(M_0, M_1)$}\\[2pt]}
\newcommand{\EncCode}[1]{\underline{\textrm{proc} $\Enc(#1)$}\\[2pt]}
\newcommand{\VfCode}[1]{\underline{\textrm{proc} $\Vf(#1)$}\\[2pt]}

\newcommand{\calK}{{\cal K}}
\newcommand{\calM}{{\cal M}}
\newcommand{\calE}{{\cal E}}
\newcommand{\calD}{{\cal D}}
\newcommand{\calN}{{\cal N}}
\newcommand{\calC}{{\cal C}}
\newcommand{\calS}{{\cal S}}
\newcommand{\calT}{{\cal T}}
\newcommand{\calH}{{\cal H}}

\newcommand{\AES}{\mathsf{AES}}
\newcommand{\GHASH}{\mathsf{GHASH}}
\newcommand{\GF}{\mathrm{GF}}
\newcommand{\CBCMAC}{\mathsf{CBC\mbox{-}MAC}}

%%%%%%%  Author Notes %%%%%%%
\ifnum\shownotes=1
\newcommand{\authnote}[2]{{ $\ll$\textsf{\footnotesize #2}$\gg$}}
\else
\newcommand{\authnote}[2]{}
\fi
%%%%%%%%%%%%%%%%%%%%%%%%%%%%%%%%%

% end of macros
%%%%%%%%%%%%%%%%%%%%%%%%%%%%%%%%%%%%%%%%%%%%%%%%%%%%%%%%%%%%%%


\title{COMPSCI 466: Homework 6}
\date{}

\usepackage{lastpage}
\pagestyle{fancy}
\lhead{\footnotesize \parbox{11cm}{COMPSCI 466, UMass Amherst, Spring 2026} }
\rhead{\footnotesize Instructor:  Adam O'Neill \\ \url{adamo@cs.umass.edu}}
\renewcommand{\headheight}{24pt}

\begin{document}

\maketitle
\thispagestyle{fancy}


\noindent
\textbf{Due date:}  May 5 11:59pm.

\bigskip

%%%%%%%%%%%%%%%%%%%%%%%%%%%%%%%%%%%%%%%%%%%%%%%%%%%%%%%%%%%%%%
\paragraph*{Problem 1. (100 points.)}
Define the symmetric encryption scheme $\schemefont{SE} = (\calK, \calE, \calD)$
where $\calK$ returns a uniformly random 128-bit key $K$, and the encryption
and decryption algorithms are as follows.

\begin{center}
\begin{tabular}{c|c}
\begin{minipage}[t]{2.1in}
\begin{tabbing}
123\=123\=123\=\kill
\textbf{Algorithm} $\calE_K(M)$: \\
\> If $|M| \neq 512$ then \return $\bot$ \\
\> Parse $M$ as $M[1] \| M[2] \| M[3] \| M[4]$ \\
\> $C_e[0] \getsr \bits^{128}$ \\
\> $C_m[0] \gets 0^{128}$ \\
\> For $i = 1, 2, 3, 4$ do: \\
\>\> $C_e[i] \gets \AES_K(C_e[i-1] \xor M[i])$ \\
\>\> $C_m[i] \gets \AES_K(C_m[i-1] \xor M[i])$ \\
\> $C_e \gets C_e[0] \| C_e[1] \| C_e[2] \| C_e[3] \| C_e[4]$ \\
\> $T \gets C_m[4]$ \\
\> \return $(C_e, T)$
\end{tabbing}
\end{minipage}
&
\begin{minipage}[t]{2.2in}
\begin{tabbing}
123\=123\=123\=\kill
\textbf{Algorithm} $\calD_K((C_e, T))$: \\
\> If $|C_e| \neq 640$ then \return $\bot$ \\
\> Parse $C_e$ as $C_e[0]\|C_e[1]\|\cdots\|C_e[4]$ \\
\> $C_m[0] \gets 0^{128}$ \\
\> For $i = 1, 2, 3, 4$ do: \\
\>\> $M[i] \gets \AES^{-1}_K(C_e[i]) \xor C_e[i-1]$ \\
\>\> $C_m[i] \gets \AES_K(C_m[i-1] \xor M[i])$ \\
\> If $C_m[4] = T$ then \return $M[1]\|M[2]\|M[3]\|M[4]$ \\
\> Else \return $\bot$
\end{tabbing}
\end{minipage}
\end{tabular}
\end{center}

\bigskip

\noindent
\textbf{(Part A --- 40 points.)}
Show that $\schemefont{SE}$ is \emph{not} IND-CPA secure.
State your adversary $A$ in pseudocode and analyze its advantage and resource usage.
\medskip

\noindent
\textbf{(Part B --- 40 points.)}
Show that $\schemefont{SE}$ is \emph{not} INT-CTXT secure.
State your adversary $A$ in pseudocode and analyze its advantage and resource usage. 

\medskip

\noindent
\textbf{(Part C --- 20 points.)} Is $\schemefont{SE}$ an Encrypt-and-MAC scheme? Briefly justify your answer.

\bigskip

%%%%%%%%%%%%%%%%%%%%%%%%%%%%%%%%%%%%%%%%%%%%%%%%%%%%%%%%%%%%%%
\paragraph*{Bonus Problem. (40 points.)}
Consider the modified scheme $\schemefont{SE}' = (\calK', \calE', \calD')$ where
$\calK'$ samples $K_e, K_m \getsr \bits^{128}$ and returns $(K_e, K_m)$.
The encryption and decryption algorithms are:

\begin{center}
\begin{tabular}{c|c}
\begin{minipage}[t]{2.2in}
\begin{tabbing}
123\=123\=123\=\kill
\textbf{Algorithm} $\calE'_{(K_e, K_m)}(M)$: \\
\> If $|M| \neq 512$ then \return $\bot$ \\
\> Parse $M$ as $M[1] \| M[2] \| M[3] \| M[4]$ \\
\> $C_e[0] \getsr \bits^{128}$ \\
\> $C_m[0] \gets 0^{128}$ \\
\> For $i = 1, 2, 3, 4$ do: \\
\>\> $C_e[i] \gets \AES_{K_e}(C_e[i-1] \xor M[i])$ \\
\>\> $C_m[i] \gets \AES_{K_m}(C_m[i-1] \xor M[i])$ \\
\> $C_e \gets C_e[0] \| C_e[1] \| C_e[2] \| C_e[3] \| C_e[4]$ \\
\> $T \gets C_m[4]$ \\
\> \return $(C_e, T)$
\end{tabbing}
\end{minipage}
&
\begin{minipage}[t]{2.3in}
\begin{tabbing}
123\=123\=123\=\kill
\textbf{Algorithm} $\calD'_{(K_e, K_m)}((C_e, T))$: \\
\> If $|C_e| \neq 640$ then \return $\bot$ \\
\> Parse $C_e$ as $C_e[0]\|C_e[1]\|\cdots\|C_e[4]$ \\
\> $C_m[0] \gets 0^{128}$ \\
\> For $i = 1, 2, 3, 4$ do: \\
\>\> $M[i] \gets \AES^{-1}_{K_e}(C_e[i]) \xor C_e[i-1]$ \\
\>\> $C_m[i] \gets \AES_{K_m}(C_m[i-1] \xor M[i])$ \\
\> If $C_m[4] = T$ then \return $M[1]\|M[2]\|M[3]\|M[4]$ \\
\> Else \return $\bot$
\end{tabbing}
\end{minipage}
\end{tabular}
\end{center}

\medskip

\noindent
\textbf{(Part A --- 20 points.)}
Is $\schemefont{SE}'$ IND-CPA secure? If yes, give convincing intuition. If no, give an attack.

\medskip

\noindent
\textbf{(Part B --- 20 points.)}
Is $\schemefont{SE}'$ INT-CTXT secure? If yes, give convincing intuition. If no, give an attack.

\end{document}
